La bioconversión de residuos orgánicos mediante larvas de mosca soldado negra (BSF, por sus siglas en inglés, \textit{Hermetia illucens}) se ha consolidado como una alternativa relevante para la valorización de biomasa residual en esquemas de economía circular. Esta relevancia se debe a su capacidad para transformar sustratos orgánicos heterogéneos en productos de valor agregado, principalmente biomasa larval rica en proteínas y lípidos, junto con un residuo transformado (\textit{frass}) de alto potencial agronómico \parencite{surendraBioconversionOrganicWastes2016,bermudezComprehensiveUtilizationBlack2023,singhInclusiveApproachOrganic2019}. 

Frente a las tecnologías convencionales de tratamiento, este proceso combina la reducción de masa residual con la recuperación de nutrientes y la generación de biomasa útil en tiempos relativamente cortos. Tales ventajas han favorecido su creciente importancia en las discusiones sobre bioeconomía, seguridad alimentaria y sostenibilidad ambiental \parencite{surendraBioconversionOrganicWastes2016,singhInclusiveApproachOrganic2019}. Esta eficiencia se sustenta en atributos biológicos distintivos de la especie, como altas tasas de ingestión, rápido crecimiento y una notable plasticidad frente a sustratos de composición variable. Estos atributos permiten convertir residuos agroindustriales, alimentarios y orgánicos urbanos en insumos para alimentación animal, extracción de compuestos industriales y sistemas agrícolas. En este sentido, la bioconversión con \textit{Hermetia illucens} no debe entenderse únicamente como una tecnología de eliminación de residuos, sino como una plataforma de transformación material y energética \parencite{bermudezComprehensiveUtilizationBlack2023,elsayedConversionProteinrichWaste2024,brunoValorizationOrganicWaste2025,surendraBioconversionOrganicWastes2016}.

En contextos agroindustriales del sur global, la acumulación y el manejo inadecuado de residuos constituyen un problema ambiental crítico. Esto se debe a su alta carga biodegradable, su generación dispersa y estacional, y una heterogeneidad fisicoquímica que dificulta su tratamiento. En regiones como el Valle del Cauca, donde confluyen actividades agrícolas y agroindustriales intensivas, esta combinación favorece pérdidas de biomasa aprovechable y aumenta los desafíos técnicos asociados al control de degradación y estabilización \parencite{foodandagricultureorganizationoftheunitednationsfaoAgriculturalWasteManagement2021,ait-kaddourTransformingPlantbasedWaste2024,nasereldeenBioconversionDifferentWaste2023}. 



En este escenario, la bioconversión mediante \textit{Hermetia illucens} adquiere interés como estrategia de valorización, precisamente por su capacidad de procesar residuos diversos en ciclos cortos. Sin embargo, su desempeño biológico y ambiental no es uniforme, sino altamente dependiente de la composición del sustrato \parencite{nayakHermetiaIllucensDiptera2024,elsayedConversionProteinrichWaste2024,brunoValorizationOrganicWaste2025}. Variaciones en la humedad, el contenido de fibra, la relación carbono/nitrógeno (C/N) y la fracción fácilmente biodegradable condicionan tanto la conversión del residuo como la dinámica metabólica del sistema. Por lo tanto, la extrapolación directa de resultados obtenidos con dietas estandarizadas de laboratorio debe hacerse con cautela \parencite{elsayedConversionProteinrichWaste2024,brunoValorizationOrganicWaste2025,rossiEstimatingDynamicsGreenhouse2024}. En consecuencia, la diversidad del sustrato no debe tratarse como un factor secundario, sino como una variable ambiental y operativa estructurante de la bioconversión.

Desde esta perspectiva, el sustrato deja de ser un simple soporte alimenticio y pasa a constituir una variable determinante del comportamiento biológico del sistema. La aptitud de un residuo para la bioconversión depende de múltiples factores, entre ellos la digestibilidad, el contenido de agua, la porosidad, el balance nutricional y la presencia de compuestos inhibitorios o carga microbiológica. En términos generales, los residuos alimentarios y subproductos agroindustriales muestran buena capacidad de conversión. Por el contrario, fracciones altamente lignocelulósicas o lodos orgánicos pueden requerir pretratamientos o presentar restricciones sanitarias y operativas \parencite{elsayedConversionProteinrichWaste2024,brunoValorizationOrganicWaste2025,nasereldeenBioconversionDifferentWaste2023}. De ello se desprende que la pregunta relevante no es solo si un residuo puede ser tratado con \textit{Hermetia illucens}, sino bajo qué condiciones fisicoquímicas la bioconversión mantiene eficiencia biológica, estabilidad y un desempeño ambiental aceptable.

En la literatura, el desempeño del proceso se ha evaluado con frecuencia mediante indicadores como la reducción de masa, el tiempo de desarrollo larval, la supervivencia y la composición proximal de la biomasa \parencite{surendraBioconversionOrganicWastes2016,singhInclusiveApproachOrganic2019,bermudezComprehensiveUtilizationBlack2023,amrulReviewOrganicWaste2022}. Aunque estos indicadores son indispensables, resultan insuficientes para describir la dinámica interna del sistema. Tienden a capturar estados finales o balances agregados, pero ofrecen poca información sobre la evolución temporal de la actividad metabólica o la respiración del sistema. Asimismo, omiten detalles sobre la competencia por oxígeno o la transición entre condiciones aeróbicas y episodios localizados de anaerobiosis \parencite{rossiEstimatingDynamicsGreenhouse2024,boakye-yiadomGreenhouseGasEmissions2022,eriksenDynamicModellingFeed2022}. Esta limitación es crítica cuando el objetivo es evaluar rigurosamente el comportamiento ambiental del proceso.

Un aspecto central del estado del arte es que la bioconversión con \textit{Hermetia illucens} suele reportarse como una alternativa favorable frente a tratamientos biológicos convencionales. Esto se atribuye a su rapidez y a su potencial para reducir la permanencia del residuo en condiciones que favorezcan fermentaciones no controladas \parencite{jenkinsProcessingPoultryManure2025,ermolaevGreenhouseGasEmissions2019,lindbergProcessEfficiencyGreenhouse2022}. Sin embargo, esta ventaja no es universal ni automática. La evidencia muestra que el desempeño depende de la calidad del sustrato, la densidad larval, la humedad, la temperatura y la disponibilidad de oxígeno. Cuando estas variables se desbalancean, el proceso puede volverse metabólicamente ineficiente y favorecer rutas microbianas no deseadas \parencite{bekkerImpactSubstrateMoisture2021,chenEffectMoistureContent2019,rossiEstimatingDynamicsGreenhouse2024,laganaroGrowthMetabolicPerformance2021}. Por ello, el análisis contemporáneo exige avanzar hacia evaluaciones situadas, condicionales y cuantificables.

En este marco, resulta relevante la transición hacia enfoques bioenergéticos y mecanísticos. Estudios recientes indican que la producción de \ce{CO2} no es solo una emisión ambiental, sino una señal fisiológica asociada a la respiración de mantenimiento y al crecimiento estructural \parencite{eriksenDynamicModellingFeed2022}. En modelos como el de Eriksen, la larva se representa mediante compartimentos fisiológicos diferenciados. Esto permite desagregar el metabolismo en costos de mantenimiento, síntesis de biomasa y acumulación de lípidos, conectando eficiencia y respiración \parencite{eriksenDynamicModellingFeed2022,grausaMetabolicModelingHermetia2023}. 

No obstante, estos modelos presentan límites en sistemas reales. Suelen centrarse en la fisiología del insecto bajo condiciones controladas, mientras que la práctica involucra la actividad de la microbiota del sustrato y la aparición de microambientes con distinta disponibilidad de oxígeno \parencite{rossiEstimatingDynamicsGreenhouse2024,bekkerImpactSubstrateMoisture2021,chenEffectMoistureContent2019}. En consecuencia, persiste una brecha entre el conocimiento del crecimiento larval como fenómeno bioenergético y la comprensión del sistema de bioconversión como un proceso ambientalmente dinámico.

En síntesis, la bioconversión con \textit{Hermetia illucens} es una tecnología de alto potencial, pero su desempeño no puede medirse solo por resultados productivos finales. La heterogeneidad de los sustratos y la interacción larva--microbiota obligan a considerar la bioconversión como un sistema biológico complejo. Su sostenibilidad depende tanto de la conversión material como de su dinámica metabólica \parencite{rossiEstimatingDynamicsGreenhouse2024,boakye-yiadomGreenhouseGasEmissions2022,eriksenDynamicModellingFeed2022}. Esta constatación conduce a examinar, en la siguiente sección, el problema específico de las emisiones de gases de efecto invernadero.

